\documentclass[10pt,a4paper]{amsart}
\usepackage[margin=19mm]{geometry}
\usepackage{amsmath,amssymb,mathtools}
\usepackage[scr=rsfs,cal=euler]{mathalfa}
\usepackage{xfrac}
\usepackage[unicode,hidelinks,bookmarks=false]{hyperref}
\newcommand{\ie}{i.e.\ }
\newcommand{\cf}{cf.\ }
\newcommand{\resp}{respectively\ }
\newcommand{\Subsection}{\subsection}
\providecommand{\nobreakdash}{\nobreak\hbox{-}}
\setlength{\emergencystretch}{2em}
\multlinegap0pt
\usepackage{bm}
\usepackage{bbm}
\usepackage{dsfont}
\usepackage{xspace}
\usepackage{cancel}
\usepackage{mathtools}
\usepackage{amsthm}
\makeatletter
\@ifundefined{theorem}{\newtheorem{theorem}{Theorem}}{}
\@ifundefined{corollary}{\newtheorem{corollary}[theorem]{Corollary}}{}
\makeatother
\makeatletter
\DeclareMathOperator{\BF}{BF}
\newcommand{\Lp}{{\mathscr{L}_p}}
\DeclareMathOperator{\bal}{bal}
\newcommand{\bb}{\mathbb}
\newcommand{\cl}{\mathcal}
\expandafter\ifx\csname example\endcsname\relax
\newenvironment{example}{\begin{examplewr}\begin{upshape}}{\end{upshape}\end{examplewr}}
\fi
\newcommand{\hf}{{\mathbf{f}}}
\newcommand{\hg}{{\mathbf{g}}}
\newcommand{\hh}{{\mathbf h}}
\expandafter\ifx\csname obs\endcsname\relax
\newenvironment{obs}{\begin{obswr}\begin{upshape}}{\end{upshape}\end{obswr}}
\fi
\expandafter\ifx\csname remark\endcsname\relax
\newenvironment{remark}{\begin{remarkwr}\begin{upshape}}{\end{upshape}\end{remarkwr}}
\fi
\newcommand{\res}{\mathrm{res}}
\makeatother
\title[Anticyclotomic diagonal classes and Beilinson--Flach element]{Anticyclotomic diagonal classes and Beilinson--Flach elements}
\author{Ra\'ul Alonso, Lois Omil-Pazos, \'Oscar Rivero}
\date{Source: arXiv:2509.07564v1 (CC BY 4.0)}
\begin{document}
\maketitle

\noindent\emph{Source and licence.} Anticyclotomic diagonal classes and Beilinson--Flach elements. Version arXiv:2509.07564v1 (CC BY 4.0), distributed under the Creative Commons Attribution 4.0 International licence, as recorded in the arXiv licence metadata for this exact version and shown on the abstract page https://arxiv.org/abs/2509.07564v1. Full rights notice: licenses/arxiv-2509.07564-CC-BY-4.0.txt.\par\smallskip

\noindent\textit{Editorial scope.} Counted unit: the Corollary whose source label is \texttt{None} in the article (source0.tex lines 1894--1899, source label \texttt{None}), delivered together with the article's own complete original proof of it (source0.tex lines 1900--1928, one \texttt{proof} environment of 29 lines). No mathematics is written, repaired or re-derived: statement and proof are the article's own text line for line. Required context retained uncounted: reclaw (lines 1109--1122); thm:reclawdiagonalcycles (lines 1278--1293). Every definition, display and label of those lines is retained unchanged. Omitted from this package and not redistributed: 1--1108, 1123--1277, 1294--1893, 1929--2620. Nothing inside a retained excerpt is omitted or altered: every retained body is delivered exactly as the article's lines are written, so no omission receipt of any kind accompanies this package. Citations: the retained excerpts carry no citation command at all, so no bibliography entry of the article is transcribed and no reference is redistributed. Reference adaptation: none was needed and none was made; every reference inside the delivered unit resolves inside it, so no unresolved reference or citation is produced. Printed slips kept literal: no printed word, symbol, hypothesis or number of the counted statement or of its proof was altered. Layout and class: the article's own class is \texttt{amsart} with packages including geometry, amsthm, amssymb, amsmath, amsfonts, mathrsfs, amscd, yfonts, verbatim, turnstile, inputenc, xy, latexsym, stmaryrd; the wrapper is the lane's pinned \texttt{amsart} editorial wrapper at 10pt on A4, which restates the theorem-like environments the retained text needs and adds the article's own macro definitions that the retained text needs (\texttt{\textbackslash BF}, \texttt{\textbackslash Lp}, \texttt{\textbackslash bal}, \texttt{\textbackslash bb}, \texttt{\textbackslash cl}, \texttt{\textbackslash hf}, \texttt{\textbackslash hg}, \texttt{\textbackslash hh}, \texttt{\textbackslash res}), with extra wrapper packages (bm, bbm, dsfont, xspace, cancel, mathtools). The delivered numbering is the unit's own. Assets: no retained span contains a figure environment, a graphics-inclusion command, a PDF-page inclusion, a table or a TikZ picture; no image file is copied into this package, the package declares no asset dependency and no image file is redistributed with it. Licence: version arXiv:2509.07564v1 of this article is distributed under the Creative Commons Attribution 4.0 International licence, as recorded in the arXiv licence metadata for this exact version and shown on the abstract page https://arxiv.org/abs/2509.07564v1. A full rights notice is delivered at licenses/arxiv-2509.07564-CC-BY-4.0.txt. Characters: every retained excerpt is pure ASCII, so the delivered engine, XeLaTeX with its default Latin Modern text font, reports no missing character.
\begin{theorem}\label{reclaw}
Fix an integer $c>1$ relatively prime to $6pN_gN_h$. Then, there exists a global cohomology class \[ _c\kappa(\hg,\hh,\mathbf{s}) \in H^1_{\bal}(\mathbb Q, \mathbb V_{\hg} \hat \otimes \mathbb V_{\hh} \hat \otimes \Lambda(1-\mathbf{s})) \] such that
\[
\cl{L}_{\hg\hh\mathbf{s}}^{\hg}(_c\kappa(\hg,\hh,\mathbf{s}))=\frac{(-1)^{\mathbf{s}}}{\lambda_{N_g}(\hg)} \cdot (c^2-c^{2\mathbf{s}-\mathbf{l}-\mathbf{m}+2}\chi_{\hg}^{(p)}(c)^{-1}\chi_{\hh}^{(p)}(c)^{-1}) \times L_p(\hg,\hh,\mathbf{s})
\]
and
\[
\cl{L}_{\hg\hh\mathbf{s}}^{\hh}(_c\kappa(\hg,\hh,\mathbf{s}))=\frac{(-1)^{\mathbf{s}}}{\lambda_{N_h}(\hh)} \cdot (c^2-c^{2\mathbf{s}-\mathbf{l}-\mathbf{m}+2}\chi_{\hg}^{(p)}(c)^{-1}\chi_{\hh}^{(p)}(c)^{-1}) \times L_p(\hh,\hg,\mathbf{s}).
\]
%\begin{enumerate}
%\item The image of the local class $\res_p(_c \kappa(\hg,\hh,\mathbf{s}))$ in $H^1(\mathbb Q_p, \mathbb V_{\hg}^- \hat \otimes \mathbb V_{\hh}\hat{\otimes}\Lambda(1-\mathbf{s}))$ lands in \[ H^1(\mathbb Q_p, \mathbb V_{\hg}^- \hat \otimes \mathbb V_{\hh}^+\hat{\otimes}\Lambda(1-\mathbf{s})). \]
%\item Letting $_c \kappa_p^{-+}(\hg,\hh)$ denote the local cohomology class in the above space, we have \[ \mathcal L_{{\bf gh}}^{-+}(_c \kappa_p^{-+}(\hg,\hh,\mathbf{s})) = \frac{(-1)^{\mathbf{s}}}{\lambda_{N_g}(\hg)} \cdot (c^2-c^{2\mathbf{s}-\mathbf{l}-\mathbf{m}+2}\chi_{\hg}^{-1}(c)\chi_{\hh}^{-1}(c)) \times L_p(\hg,\hh,\mathbf{s}). \] %where $\lambda_{\hg}$ denotes the pseudo-eigenvalue of $\hg$, an Iwasawa function interpolating the pseudo-eigenvalue at $N$ of the crystalline classical specializations of $\hg$.
%\end{enumerate}
\end{theorem}

\begin{theorem}\label{thm:reclawdiagonalcycles}
There exists a global cohomology class
\[
\kappa(\hf,\hg,\hh) \in H^1_{\bal}(\mathbb Q, \mathbb V_{\hf} \hat \otimes \mathbb V_{\hg} \hat \otimes \mathbb V_{\hh}(2-{\bf t}))
\]
such that, for $\bm{\xi}\in \lbrace \hf,\hg,\hh\rbrace$,
we have that
\[
\mathcal L_{\hf\hg\hh}^{\bm{\xi}}(\kappa(\hf, \hg, \hh)) = \Lp^{\xi}(\hf, \hg, \hh).
\]
%\begin{enumerate}
%\item The projection of the local class $\res_p(\kappa(\hf,\hg,\hh))$ to $H^1(\mathbb Q_p, \mathbb V_{\hf}^- \hat \otimes \mathbb V_{\hg} \hat \otimes \mathbb V_{\hh}(1-{\bf c}))$ lands in $H^1(\mathbb Q_p, \mathbb V_{\hf}^- \hat \otimes \mathbb V_{\hg}^+ \hat \otimes \mathbb V_{\hh}^+(1-{\bf c}))$.
%\item Letting $\kappa_p^{-++}(\hf,\hg,\hh)$ denote the local cohomology class in the above space, \[ \mathcal L_{\hf \hg \hh}^{-++}(\kappa_p^{-++} (\hf, \hg, \hh)) = \Lp^f(\hf, \hg, \hh). \]
%\end{enumerate}

\end{theorem}

\begin{corollary}
    Assume that $H^1_{\cl{G}\cup +}(\bb{Q}, \bb{V}_{f\hg\hh}^\dagger)$ is a torsion-free $\Lambda_{\hg\hh}$-module of rank $1$. Then
    \begin{align*}
     \Lp^g(f,\hg,\hh)\Lp^h(f,\hg,\hh)&=\frac{\lambda_{N_g}(\hg)}{\lambda_{N_h}(\hh)}\cdot L_p\left(\hg,\hh,\frac{\mathbf{l}+\mathbf{m}-1}{2}\right)L_p\left(\hh\otimes\varepsilon_K,\hg,\frac{\mathbf{l}+\mathbf{m}-1}{2}\right) \\ &=\frac{\lambda_{N_g}(\hg)}{\lambda_{N_h}(\hh)}\cdot L_p\left(\hg,\hh\otimes\varepsilon_K,\frac{\mathbf{l}+\mathbf{m}-1}{2}\right)L_p\left(\hh,\hg,\frac{\mathbf{l}+\mathbf{m}-1}{2}\right).
    \end{align*}
\end{corollary}

\begin{proof}
    By the previous corollary, we know that $\BF(f,\hg,\hh)\in H^1_{\bal}(\bb{Q},\bb{V}_{f\hg\hh}^\dagger)$. Therefore, the image of $\BF(f,\hg,\hh)$ in $H^1(\bb{Q}_p,V_f^-\otimes\bb{V}_{\hg}\hat{\otimes}\bb{V}_{\hh}^{-}(2-\mathbf{t}_1))$ is zero. This image is given by
    \[
    v_f^-\otimes \left(L_p\left(\hg,\hh\otimes\varepsilon_K,\frac{\mathbf{l}+\mathbf{m}-1}{2}\right)\res_p(\kappa_{\hg,\hh})^{+-}-L_p\left(\hg,\hh,\frac{\mathbf{l}+\mathbf{m}-1}{2}\right)\res_p(\kappa_{\hg,\hh\otimes\varepsilon_K})^{+-}\right),
    \]
    so it follows that
    \[
    L_p\left(\hg,\hh\otimes\varepsilon_K,\frac{\mathbf{l}+\mathbf{m}-1}{2}\right)\res_p(\kappa_{\hg,\hh})^{+-}-L_p\left(\hg,\hh,\frac{\mathbf{l}+\mathbf{m}-1}{2}\right)\res_p(\kappa_{\hg,\hh\otimes\varepsilon_K})^{+-}=0.
    \]
    Applying the Perrin-Riou map $\mathcal{L}_{\hg\hh}^{+-}$, we deduce by Theorem~\ref{reclaw} that
    \begin{align}
    &L_p\left(\hg,\hh,\frac{\mathbf{l}+\mathbf{m}-1}{2}\right)L_p\left(\hh\otimes\varepsilon_K,\hg,\frac{\mathbf{l}+\mathbf{m}-1}{2}\right) \nonumber \\
    &=L_p\left(\hg,\hh\otimes\varepsilon_K,\frac{\mathbf{l}+\mathbf{m}-1}{2}\right)L_p\left(\hh,\hg,\frac{\mathbf{l}+\mathbf{m}-1}{2}\right). \label{eq:equalityofHRpadicLfunctions}
    \end{align}
    Now note that, for any element
    \[
    v\otimes z \in V_f^+ \otimes H^1(\bb{Q}_p,\bb{V}_\hg^+\hat{\otimes} \bb{V}_\hh^-(2-\mathbf{t}_1))\cong H^1(\bb{Q}_p,V_f^+\otimes \bb{V}_{\hg}^+\hat{\otimes}\bb{V}_{\hh}^-(2-\mathbf{t}_1)),
    \]
    we have that
    \[
    \mathcal{L}_{f\hg\hh}^{++-}(v\otimes z)=\langle v,\omega_f\rangle \cdot\mathcal{L}_{\hg\hh}^{+-}(z).
    \]
    Let $\mathrm{pr}^{++-}:H^1_{\bal}(\bb{Q}_p,\bb{V}_{f\hg\hh}^\dagger)\rightarrow H^1(\bb{Q}_p,V_f^+\otimes\bb{V}_\hg^+\hat{\otimes}\bb{V}_\hh^-(2-\mathbf{t}_1))$ be the map induced by the natural projection $\mathscr{F}^2\bb{V}_{f\hg\hh}^\dagger\rightarrow V_f^+\otimes\bb{V}_\hg^+\hat{\otimes}\bb{V}_\hh^-(2-\mathbf{t}_1)$. Using Theorem~\ref{reclaw} and equation~\ref{eq:equalityofHRpadicLfunctions}, the image of $\BF(f,\hg,\hh)$ by $\cl{L}_{f\hg\hh}^{++-}\circ \mathrm{pr}^{++-}\circ \res_p$ is equal to
    \[
    \Omega_{f,\gamma}\cdot\lambda_{N_h}(\hh)^{-1}\cdot L_p\left(\hg,\hh,\frac{\mathbf{l}+\mathbf{m}-1}{2}\right)L_p\left(\hh\otimes\varepsilon_K,\hg,\frac{\mathbf{l}+\mathbf{m}-1}{2}\right).
    \]
    Also, by Theorem~\ref{thm:reclawdiagonalcycles}, the image of $\kappa(f,\hg,\hh)$ by $\cl{L}_{f\hg\hh}^{++-}\circ \mathrm{pr}^{++-}\circ \res_p$ is equal to $\Lp^h(f,\hg,\hh)$. Hence, the result now follows from the previous corollary.
    
\end{proof}

\end{document}
